\chapter*{Conclusion générale}
\addcontentsline{toc}{chapter}{Conclusion générale}
Le problème Ring-Star fournit un cadre clair pour étudier ensemble le choix
des stations et leur ordre de desserte sur une ligne circulaire. La
formulation retenue fixe la cardinalité, impose une station et minimise
une somme pondérée de longueur d'anneau et de distances d'accès.
L'intégration du voyageur de commerce comme cas particulier explique sa
NP-difficulté et le besoin de méthodes complémentaires.

La PLNE avec flot exprime directement la sélection, les affectations et
la connexité du cycle. Elle apporte une garantie lorsque CBC termine avec
une preuve, ou une solution réalisable accompagnée d'une borne lorsqu'il
s'arrête avant. Les méthodes approchées ont un rôle différent : construire
rapidement une solution, puis explorer des échanges de stations en
réévaluant les deux composantes de l'objectif.

Sur les 24 cas de la campagne principale, les recherches améliorent les
coûts moyens de la construction. Le gain médian atteint 11,79\,\% pour
la locale et 13,67\,\% pour le tabou. La meilleure moyenne tabou dans
22 cas s'accompagne d'un effort de calcul plus important. Le résultat
doit donc être présenté comme un compromis observé sous des réglages
précis, et non comme une supériorité démontrée à ressources égales.

La comparaison à un optimum est possible sur trois petits cas. Sur
Ulysses16 à cinq stations, les trois trajectoires tabou atteignent le
coût prouvé de 65,3775. Sur les grandes instances, les solutions
restent réalisables et améliorées, sans preuve de leur proximité à
l'optimum. L'exemple att48 illustre cette distinction : une baisse de
12,47\,\% est mesurée par rapport à la construction, avec un coût final
de 42\,743,75, mais sans certificat d'optimalité.

L'organisation des sorties fait partie du résultat du projet. Les
solutions, paramètres, traces et versions permettent de retrouver
l'origine d'un chiffre et de redessiner les figures. La relecture
indépendante des coûts et les tests complètent cette démarche.
Le défaut de dernière ligne ignorée repéré dans le CSV rappelle
l'intérêt de rapprocher les fichiers entre eux, même lorsque les
solutions et les synthèses sont correctes.

La suite du travail peut se concentrer sur des budgets équitables,
davantage de graines et une étude du poids $\alpha$. Une application
opérationnelle demanderait en parallèle de remplacer les distances
planes et les poids uniformes par des données de déplacement,
des demandes et des contraintes réelles. Le projet constitue ainsi
une base expérimentale documentée pour explorer la localisation-routage,
avec des résultats dont la portée et les limites sont explicites.
